\changecolours{ScoutGreen}
\section[Filas e Desacoplamento]{Filas Assíncronas de Checkout com Listas do Redis}

% ------------------------------------------------------------------------------
\twocolframe[title={O Problema do Checkout Síncrono no E-commerce},leftcol=7.1cm,rightcol=6.9cm]{%
  \alert{Por que o Checkout Síncrono Falha?}\par
  \itemseps[topsepi=0.2mm,itemsepi=0.4mm,labelsep=1.5mm,leftmargini=2.5mm]
  \begin{itemize}
    \item \textbf{Cadeia Lenta:} Cartão, antifraude, NF-e e e-mail levam de 2 a 6 segundos.
    \item \textbf{Esgotamento:} 1000 clientes simultâneos derrubam servidores por conexões abertas.
    \item \textbf{Solução com Fila:} A API valida o carrinho, empacota e joga na fila do Redis ($\approx 0.5\text{ ms}$).
  \end{itemize}
}{%
  \alert{O Padrão Produtor-Consumidor FIFO}\par
  \itemseps[topsepi=0.2mm,itemsepi=0.4mm,labelsep=1.5mm,leftmargini=2.5mm]
  \begin{itemize}
    \item \textbf{Produtor (API Web):}
      Insere na esquerda com \texttt{LPUSH fila:pedidos <dados>}.
    \item \textbf{Consumidor (Worker):}
      Retira na direita com \texttt{BRPOP fila:pedidos 5}.
    \item \textbf{Por que \texttt{BRPOP}?}\par
      Não consome CPU em loop (\textit{busy waiting}). O processo dorme no kernel até um item chegar!
  \end{itemize}
}

% ------------------------------------------------------------------------------
\begin{frame}[fragile]{Implementação do Produtor e do Consumidor}
  \vspace{0.1cm}
  \begin{columns}[t]
    \begin{column}{0.48\textwidth}
      \footnotesize
      \alert{Produtor (API Web no Checkout)}\par
      \vspace{0.1cm}
\begin{lstlisting}[language=Python,basicstyle=\fontsize{4.6pt}{5.6pt}\selectfont\ttfamily]
@app.route("/api/checkout", methods=["POST"])
def checkout():
    # 1. Le carrinho do Redis Hash
    itens = r.hgetall(f"cart:{session_id}")
    if not itens:
        return jsonify({"erro": "vazio"}), 400

    pedido = {
        "id_pedido": f"PED-{int(time.time())}",
        "session_id": session_id,
        "total": calcular_total(itens),
        "status": "ENFILEIRADO"
    }
    
    # 2. LPUSH: Enfileira pedido na lista
    pos = r.lpush("fila:pedidos", json.dumps(pedido))
    # 3. Limpa o carrinho
    r.delete(f"cart:{session_id}")
    return jsonify({"pedido": pedido, "pos": pos})
\end{lstlisting}
    \end{column}

    \begin{column}{0.48\textwidth}
      \alert{Consumidor (\texttt{worker\_pedidos.py})}\par
      \vspace{0.1cm}
\begin{lstlisting}[language=Python,basicstyle=\fontsize{4.6pt}{5.6pt}\selectfont\ttfamily]
def processar_pedidos():
    chave = "fila:pedidos"
    print("Worker escutando via BRPOP...")
    
    while True:
        # BRPOP bloqueia ate 5s sem CPU
        res = r.brpop(chave, timeout=5)
        if res is None:
            continue
            
        _, pedido_raw = res
        pedido = json.loads(pedido_raw)
        
        print(f"Processando: {pedido['id_pedido']}")
        executar_pagamento(pedido)
        
        # PREVIEW DA AULA 06 (MONGODB):
        # db.pedidos.insert_one(pedido)
\end{lstlisting}
    \end{column}
  \end{columns}
\end{frame}
